\documentclass[12pt,a4paper]{article}
\usepackage[margin=1.45cm]{geometry}
\usepackage{newtxtext, newtxmath}
\usepackage{graphicx}
\usepackage{booktabs}
\usepackage{array}
\usepackage{enumitem}
\usepackage{setspace}
\setlength{\parindent}{0pt}
\setlength{\parskip}{0.45em}
\pagestyle{empty}

\begin{document}

\begin{center}
\begin{spacing}{1.5}
{\Large\textbf{Design and Hardware Implementation of a Real-Time Image Tracking System using Verilog HDL}}\\
	\textbf{Ayush Purohit (AU2240179), BTech. Computer Science and Engineering}\\
	\textbf{Ahmedabad University}\\
	\textbf{Mentor(s):} Internal Mentor: Prof. Mazad Zaveri
\end{spacing}
\end{center}

This project adopts a hardware-friendly object tracking pipeline that is based on SORT. 
(Basic Online and Realtime Tracking). The overall goal is to localize the software. 
convert predict-update tracking cycle into a Verilog design which is deterministic, 
real-time operation. The video stream is initially processed with MOG2 in a Python front-end. 
subtraction of background, contour filtering and temporal smoothing to find the most predominant. 
object per frame. These detections are exported using various formats to ensure that both software are used. 
aligned frame-wise inputs are received by reference and hardware controller.

The tracker state can be expressed as $[c_x, c_y, h, r, v_x, v_y, v_h]^T$ with 16-bit fixed-point. 
registers (FXP(8,8)). The controller is to be a finite-state machine whose stages are: 
Idle, Predict, IoU Gate, Update/No-Update, Metadata Update, and Done. Prediction follows 
a constant-velocity model and IoU gating determine the use of measurement correction. 
or project the anticipated state. Metadata of trackers is also maintained in the implementation. 
emits status messages like matched, save, reset, done, (age, hits, time-since-update) 
and busy in order to be integrated with higher-level logic.

\begin{figure}[h]
\centering
\includegraphics[width=0.86\linewidth,height=0.25\textheight,keepaspectratio]{Specs/Figures/FlowChart.pdf}
\caption{Setup/Product Architecture: End-to-end flow for detection, prediction, IoU association, and hardware controller update.}
\end{figure}

\vspace{-0.9em}
\begin{table}[h]
\centering
\small
\begin{tabular}{ | p{0.14\linewidth} | p{0.15\linewidth} | p{0.14\linewidth} | p{0.15\linewidth} | p{0.14\linewidth} | }
\hline
\textbf{Data} & \textbf{Detection} & \textbf{Prediction} & \textbf{Association} & \textbf{Validation}\\
\hline
Video frames converted to synchronized CSV streams (HW + SORT). & OpenCV MOG2 + morphology + contour-area threshold + smoothing. & Constant-velocity Kalman-style state advance in Verilog fixed-point datapath. & IoU-gated update/no-update decision with metadata counters. & Frame-aligned comparison against Python SORT reference output.\\
\hline
\end{tabular}
\end{table}

\vspace{-0.8em}
The project shows that hardware-oriented tracking with Verilog is an end-to-end feasible project. 
but retaining the key SORT workflow. The prototype at hand is functionally. 
integrated and appropriate as a starting point towards further optimization, including detector. 
calibration, scaling/aligning of coordinates, and future multi-object/FPGA implementation activities.
\end{document}